\section{互信息是依赖，也是对不确定性的减少}
\label{sec:07-Information-Theory/02-Joint-Information/02}

两个变量，你能用 Pearson 相关系数检验线性耦合。但 $Y=X^2$ 的相关系数可以接近零——不是没关系，而是关系不是线性的。你需要一个工具，不问关系长什么样，只问有没有关系。

互信息干了这件事。

\subsection{直觉入口：知道一个后对另一个少迷惑了多少}

观察 $Y$ 之前，对 $X$ 的不确定性是 $H(X)$；观察 $Y$ 之后，平均还剩 $H(X\mid Y)$。二者之差自然衡量了 $Y$ 告诉你多少关于 $X$ 的信息：

\[
I(X;Y)=H(X)-H(X\mid Y).
\]

这个差值就是互信息。数值越大，知道 $Y$ 后关于 $X$ 的迷惑减少越多。

\subsection{它是对称的——共享的信息没有方向}

你可能觉得``$Y$ 告诉 $X$''这话带方向。但把熵链式法则展开：

\[
\begin{aligned}
H(X,Y) &= H(X)+H(Y\mid X) \\
&= H(Y)+H(X\mid Y).
\end{aligned}
\]

整理：

\[
\begin{aligned}
I(X;Y) &= H(X)-H(X\mid Y) \\
&= H(Y)-H(Y\mid X) \\
&= H(X)+H(Y)-H(X,Y).
\end{aligned}
\]

从 $H(X)-H(X\mid Y)$ 和 $H(Y)-H(Y\mid X)$ 是同一个数这件事，就知道互信息对称：

\[
I(X;Y)=I(Y;X).
\]

语言上有``X 告诉 Y''和``Y 告诉 X''的方向感，但统计共享的信息分量本身没有方向——它是两个变量共有的那部分不确定性。

\subsection{互信息就是联合分布离独立有多远}

把定义展开，不加推导直接看结果：

\[
I(X;Y)=\sum_{x,y}p(x,y)\log\frac{p(x,y)}{p(x)p(y)}.
\]

这就是联合分布 $p_{XY}$ 与``假装独立''的乘积分布 $p_X p_Y$ 之间的 KL 散度：

\[
I(X;Y)=D(p_{XY}\|p_X p_Y),
\]

其中 $D(p\|q)=\sum_x p(x)\log\frac{p(x)}{q(x)}$。KL 散度非负且为零当且仅当两分布相同（下一章会系统处理它），因此

\[
I(X;Y)\ge 0,
\]

且

\[
I(X;Y)=0\iff\text{$X$ 与 $Y$ 独立}.
\]

这一结果在实践中分量很重。Pearson 系数为零只排除线性耦合，互信息为零等价于完全独立——任何形式的依赖，线性的、二次的、周期的、乃至任何可测的耦合，都会让互信息大于零。一个数字，检测一切依赖。

\subsection{上界和确定关系}

离散情形下条件熵非负，所以

\[
I(X;Y)\le H(X),\qquad I(X;Y)\le H(Y).
\]

结合两条：

\[
I(X;Y)\le\min\{H(X),H(Y)\}.
\]

互信息的上界是各自熵的较小者——一个变量携带的关于另一个变量的信息，不可能超过它自身包含的全部信息，这句话几乎是同义反复。

若 $Y=g(X)$ 是 $X$ 的确定函数，则 $H(Y\mid X)=0$，于是

\[
I(X;Y)=H(Y).
\]

$g$ 若是一一对应，$H(Y)=H(X)$——全部信息保留；若是多对一，则只有 $g(X)$ 还能区分的那些部分被保留，其余的丢了。

\subsection{二元对称信道：噪声和信息}

继续上一节的例子。输入 $X\sim\mathrm{Bernoulli}(1/2)$，噪声 $N\sim\mathrm{Bernoulli}(\varepsilon)$ 独立，输出 $Y=X\oplus N$。已算出 $H(Y)=1$，$H(Y\mid X)=h_2(\varepsilon)$，因此

\[
I(X;Y)=1-h_2(\varepsilon).
\]

看三个特征点：

\begin{itemize}
\tightlist
\item $\varepsilon=0$：无噪声，互信息为 1 bit——输出完美告诉你输入；
\item $\varepsilon=1/2$：完全随机噪声，输出与输入独立，互信息为 0；
\item $\varepsilon=1$：输出总是输入的反，但你可以把输出翻回来，互信息又是 1 bit。
\end{itemize}

第三个点最值得停顿一下：确定性翻转虽然``每位都错''，却不丢任何信息。你只要知道翻转规律，就能完全恢复输入。真正吃掉信息的是不可预测的随机噪声，不是确定性失真。

高错误率不等于低信息量——这个区分在分析含噪信道、或者比较不同传感器品质时很实用。

\subsection{连续互信息不受坐标干扰}

微分熵随坐标变化的问题，到了互信息这里被化解了。直接用测度层面的 KL 散度定义：

\[
I(X;Y)=D(P_{XY}\|P_XP_Y).
\]

当联合分布与边缘分布都有密度时，写成积分：

\[
I(X;Y)=\iint f(x,y)\log\frac{f(x,y)}{f_X(x)f_Y(y)}\,dx\,dy.
\]

对 $X$ 和 $Y$ 分别做光滑可逆变换 $U=g(X), V=h(Y)$：$h(U)$ 增加的 $\E\log|\det J_g|$ 项与 $h(V)$ 增加的 $\E\log|\det J_h|$ 项的和，恰好等于 $h(U,V)$ 增加的对应两项之和，在互信息的熵差形式中完全抵消。因此

\[
I(U;V)=I(X;Y).
\]

互信息在分别可逆的重参数化下不变。这让它成为连续域里比单独微分熵更可靠的依赖度量——你不必担心换了把尺子答案就会变。

\subsection{估计互信息是另一回事}

离散小概率表上直接点估计通常还行。高维连续数据里，你往往需要估计密度比，或者借助变分下界、kNN 等非参数方法。有限样本下的估计通常有偏；基于神经网络的估计器也可能因为优化不充分或模型容量不足而给出严重偏差。

把一个训练目标命名为``互信息下界''，不等于跑出来的数值就是真实互信息的紧致估计。估计方法和假设本身需要独立的验证——和任何统计量一样。
